%\VignetteIndexEntry{practical3}
%!Snw weave = knitr
%\VignetteEngine{knitr::knitr}

\documentclass[a4paper,justified,openany]{tufte-handout}\usepackage[]{graphicx}\usepackage[]{color}
%% maxwidth is the original width if it is less than linewidth
%% otherwise use linewidth (to make sure the graphics do not exceed the margin)
\makeatletter
\def\maxwidth{ %
  \ifdim\Gin@nat@width>\linewidth
    \linewidth
  \else
    \Gin@nat@width
  \fi
}
\makeatother

\definecolor{fgcolor}{rgb}{0.345, 0.345, 0.345}
\newcommand{\hlnum}[1]{\textcolor[rgb]{0.686,0.059,0.569}{#1}}%
\newcommand{\hlstr}[1]{\textcolor[rgb]{0.192,0.494,0.8}{#1}}%
\newcommand{\hlcom}[1]{\textcolor[rgb]{0.678,0.584,0.686}{\textit{#1}}}%
\newcommand{\hlopt}[1]{\textcolor[rgb]{0,0,0}{#1}}%
\newcommand{\hlstd}[1]{\textcolor[rgb]{0.345,0.345,0.345}{#1}}%
\newcommand{\hlkwa}[1]{\textcolor[rgb]{0.161,0.373,0.58}{\textbf{#1}}}%
\newcommand{\hlkwb}[1]{\textcolor[rgb]{0.69,0.353,0.396}{#1}}%
\newcommand{\hlkwc}[1]{\textcolor[rgb]{0.333,0.667,0.333}{#1}}%
\newcommand{\hlkwd}[1]{\textcolor[rgb]{0.737,0.353,0.396}{\textbf{#1}}}%
\let\hlipl\hlkwb

\usepackage{framed}
\makeatletter
\newenvironment{kframe}{%
 \def\at@end@of@kframe{}%
 \ifinner\ifhmode%
  \def\at@end@of@kframe{\end{minipage}}%
  \begin{minipage}{\columnwidth}%
 \fi\fi%
 \def\FrameCommand##1{\hskip\@totalleftmargin \hskip-\fboxsep
 \colorbox{shadecolor}{##1}\hskip-\fboxsep
     % There is no \\@totalrightmargin, so:
     \hskip-\linewidth \hskip-\@totalleftmargin \hskip\columnwidth}%
 \MakeFramed {\advance\hsize-\width
   \@totalleftmargin\z@ \linewidth\hsize
   \@setminipage}}%
 {\par\unskip\endMakeFramed%
 \at@end@of@kframe}
\makeatother

\definecolor{shadecolor}{rgb}{.97, .97, .97}
\definecolor{messagecolor}{rgb}{0, 0, 0}
\definecolor{warningcolor}{rgb}{1, 0, 1}
\definecolor{errorcolor}{rgb}{1, 0, 0}
\newenvironment{knitrout}{}{} % an empty environment to be redefined in TeX

\usepackage{alltt}




\usepackage{amsmath}

% Set up the images/graphics package
\usepackage{graphicx}
\setkeys{Gin}{width=\linewidth,totalheight=\textheight,keepaspectratio}
\graphicspath{{graphics/}}

\title{Building an R package}
\author[Jumping Rivers]{Jumping Rivers}
\date{}  % if the \date{} command is left out, the current date will be used

\usepackage{booktabs}
\usepackage{units}
\usepackage{fancyvrb}
\fvset{fontsize=\normalsize}
\newcommand{\cc}{\texttt}
\graphicspath{{../graphics/}}
\setcounter{secnumdepth}{2}
\usepackage{microtype}
\IfFileExists{upquote.sty}{\usepackage{upquote}}{}
\begin{document}

\maketitle% this prints the handout title, author, and date
\begin{marginfigure}
\centering
\includegraphics[]{logo.png}
\end{marginfigure}
\marginnote{Do you need statistical or computational assistance in a grant application?
We can help!}
\section{My first package}

\subsection{RStudio projects}

RStudio projects make it straightforward to switch between analyses. You can
create an RStudio project:
\begin{itemize}
\item in a brand new directory or
\item in an existing directory where you already have R code and data or
\item by cloning a version control\sidenote{Git or Subversion} repository.
\end{itemize}
Each project has its own working directory, workspace, history, and source
documents.

When we create an RStudio package project, this will create a new directory that
contains five files.
\begin{itemize}
\item \cc{.RbuildIgnore}: this file contains a list of file names that the R
  build process should ignore when constructing your package. For example, the \cc{.Rproj} file.
\item \cc{DESCRIPTION}: an overview of your package.\sidenote{\url{http://goo.gl/AXSft}}
\item \cc{*.Rproj}: an RStudio project file.
\item \cc{NAMESPACE}: a file that contains exported/exposed package functions\sidenote{\url{http://goo.gl/I5yGf8}}.
\end{itemize}
Your package directory will also contain two directories.
\begin{itemize}
\item \cc{man}: documentation directory: when you use the \cc{help} function in R, the help pages returned live in
this directory.
\item \cc{R}: R file directory: all R functions exported from your package belong in this directory.
\end{itemize}

\subsection*{Tasks}

\begin{enumerate}
\item Create  a ``package'' project, via
\begin{verbatim}
    File -> New Project -> New Directory -> R package
\end{verbatim}
Call the package \cc{pkg} and select the directory you want to store the package
in. Then click
\begin{verbatim}
    Create Project
\end{verbatim}
Congratulations - you have just created your first package called \cc{pkg}
\item Click
\begin{verbatim}
Build -> Build and Reload
\end{verbatim}
\item Now type
\begin{knitrout}
\definecolor{shadecolor}{rgb}{0.969, 0.969, 0.969}\color{fgcolor}\begin{kframe}
\begin{alltt}
\hlkwd{library}\hlstd{(}\hlstr{"pkg"}\hlstd{)}
\hlkwd{hello}\hlstd{()}
\end{alltt}
\end{kframe}
\end{knitrout}
\item Look at the \cc{NAMESPACE} file. Notice that the function \cc{hello}
  has been exported.
\item The \cc{hello} function has also been documented
\begin{knitrout}
\definecolor{shadecolor}{rgb}{0.969, 0.969, 0.969}\color{fgcolor}\begin{kframe}
\begin{alltt}
\hlkwd{help}\hlstd{(}\hlstr{"hello"}\hlstd{)}
\end{alltt}
\end{kframe}
\end{knitrout}
The associated documentation file is in the \cc{man} directory.
\end{enumerate}

\subsection{roxygen2}

Keeping the \cc{NAMESPACE} file and documentation up-to-date is a painful
experience. To ameliorate the process, we use the \cc{roxygen2} package to
automatically generate the necessary entries. Above function definitions we add
\cc{roxygen2} tags. The tags are of the form:

\begin{knitrout}
\definecolor{shadecolor}{rgb}{0.969, 0.969, 0.969}\color{fgcolor}\begin{kframe}
\begin{alltt}
\hlcom{#' @export}
\hlcom{#' @details}
\hlcom{#' @aliases}
\end{alltt}
\end{kframe}
\end{knitrout}
\noindent Notice the tags are just R comments.

\subsection*{Tasks}

\begin{enumerate}
\item Check that you have the necessary R packages installed
\begin{knitrout}
\definecolor{shadecolor}{rgb}{0.969, 0.969, 0.969}\color{fgcolor}\begin{kframe}
\begin{alltt}
\hlkwd{library}\hlstd{(}\hlstr{"devtools"}\hlstd{)}
\hlkwd{library}\hlstd{(}\hlstr{"roxygen2"}\hlstd{)}
\end{alltt}
\end{kframe}
\end{knitrout}
If you don't have them installed, then install them in the usual way
\begin{knitrout}
\definecolor{shadecolor}{rgb}{0.969, 0.969, 0.969}\color{fgcolor}\begin{kframe}
\begin{alltt}
\hlkwd{install.packages}\hlstd{(}\hlkwd{c}\hlstd{(}\hlstr{"devtools"}\hlstd{,} \hlstr{"roxygen2"}\hlstd{))}
\end{alltt}
\end{kframe}
\end{knitrout}
\item Click on
\begin{verbatim}
    Build -> Configure build tools
\end{verbatim}
then select
\begin{verbatim}
    Generate documentation with Roxygen
\end{verbatim}
and click OK. Now when we build our package, RStudio will automatically run
\begin{knitrout}
\definecolor{shadecolor}{rgb}{0.969, 0.969, 0.969}\color{fgcolor}\begin{kframe}
\begin{alltt}
\hlkwd{library}\hlstd{(}\hlstr{"roxygen2"}\hlstd{)}
\hlkwd{roxygenise}\hlstd{(}\hlstr{"."}\hlstd{)}
\end{alltt}
\end{kframe}
\end{knitrout}
\item In the file \cc{R/hello.R} add
\begin{knitrout}
\definecolor{shadecolor}{rgb}{0.969, 0.969, 0.969}\color{fgcolor}\begin{kframe}
\begin{alltt}
\hlcom{#' @export}
\end{alltt}
\end{kframe}
\end{knitrout}
just above the \cc{hello} function, i.e.
\begin{knitrout}
\definecolor{shadecolor}{rgb}{0.969, 0.969, 0.969}\color{fgcolor}\begin{kframe}
\begin{alltt}
\hlcom{#' @export}
\hlstd{hello} \hlkwb{<-} \hlkwa{function}\hlstd{() \{}
  \hlkwd{print}\hlstd{(}\hlstr{"Hello, world!"}\hlstd{)}
\hlstd{\}}
\end{alltt}
\end{kframe}
\end{knitrout}
The \cc{export} tag above the \cc{hello} function indicates that we
want to export\sidenote{Export means the users loading this package, can access
  this function.} this particular function.
\item Now delete the \cc{man} directory and the \cc{NAMESPACE} file\sidenote{The
    reason for deleting these is that we will automatically generate them using
    \cc{roxygen2} - more details below.}. Select\sidenote{The keyboard shortcut
    for this is \mbox{\cc{Ctrl + B}}}
\begin{verbatim}
    Build -> Build and reload
\end{verbatim}
You should now be able to load your package with
\begin{knitrout}
\definecolor{shadecolor}{rgb}{0.969, 0.969, 0.969}\color{fgcolor}\begin{kframe}
\begin{alltt}
\hlkwd{library}\hlstd{(}\hlstr{"pkg"}\hlstd{)}
\end{alltt}
\end{kframe}
\end{knitrout}
\noindent and call the \cc{hello} function
\begin{knitrout}
\definecolor{shadecolor}{rgb}{0.969, 0.969, 0.969}\color{fgcolor}\begin{kframe}
\begin{alltt}
\hlkwd{hello}\hlstd{()}
\end{alltt}
\end{kframe}
\end{knitrout}
\item Open the \cc{NAMESPACE} file. You should see the entry
\begin{knitrout}
\definecolor{shadecolor}{rgb}{0.969, 0.969, 0.969}\color{fgcolor}\begin{kframe}
\begin{alltt}
\hlcom{# Generated by roxygen2: do not edit by hand}
\hlkwd{export}\hlstd{(hello)}
\end{alltt}
\end{kframe}
\end{knitrout}
\end{enumerate}


\subsection{Adding new functions}


\newthought{All R} functions that we create in our package are saved in the
\cc{R/} directory\sidenote{The files can have a \cc{.R} or \cc{.r} file
  extension. Personally, I prefer \cc{.R}}. This directory can contain multiple
files.

\subsection*{Tasks}

\begin{enumerate}
\item Create a new file in the \cc{R/} directory called \cc{basic.R}. In this
  file add the following code
\begin{knitrout}
\definecolor{shadecolor}{rgb}{0.969, 0.969, 0.969}\color{fgcolor}\begin{kframe}
\begin{alltt}
\hlcom{#' @export}
\hlstd{add} \hlkwb{=} \hlkwa{function}\hlstd{(}\hlkwc{x}\hlstd{,} \hlkwc{y}\hlstd{) \{}
  \hlkwd{return}\hlstd{(x} \hlopt{+} \hlstd{y)}
\hlstd{\}}
\end{alltt}
\end{kframe}
\end{knitrout}
\noindent Build and reload your package\sidenote{Remember the keyboard shortcut
  \mbox{\cc{Ctrl + B}}}. After reloading your package, the following code should
run
\begin{knitrout}
\definecolor{shadecolor}{rgb}{0.969, 0.969, 0.969}\color{fgcolor}\begin{kframe}
\begin{alltt}
\hlkwd{library}\hlstd{(}\hlstr{"pkg"}\hlstd{)}
\hlkwd{add}\hlstd{(}\hlnum{1}\hlstd{,} \hlnum{2}\hlstd{)}
\end{alltt}
\end{kframe}
\end{knitrout}

\item Create a new function called \cc{check\_numeric}
\begin{knitrout}
\definecolor{shadecolor}{rgb}{0.969, 0.969, 0.969}\color{fgcolor}\begin{kframe}
\begin{alltt}
\hlstd{check_numeric} \hlkwb{=} \hlkwa{function}\hlstd{(}\hlkwc{x}\hlstd{)} \hlkwd{all}\hlstd{(}\hlkwd{is.numeric}\hlstd{(x))}
\end{alltt}
\end{kframe}
\end{knitrout}
\noindent and save it in the \cc{basic.R} file.
\item Now use \cc{check\_numeric} in the \cc{add} function
\begin{knitrout}
\definecolor{shadecolor}{rgb}{0.969, 0.969, 0.969}\color{fgcolor}\begin{kframe}
\begin{alltt}
\hlstd{add} \hlkwb{=} \hlkwa{function}\hlstd{(}\hlkwc{x}\hlstd{,} \hlkwc{y}\hlstd{) \{}
  \hlkwa{if}\hlstd{(}\hlopt{!}\hlkwd{check_numeric}\hlstd{(}\hlkwd{c}\hlstd{(x, y)))} \hlkwd{stop}\hlstd{(}\hlstr{"Not numeric"}\hlstd{)}
  \hlstd{x} \hlopt{+} \hlstd{y}
\hlstd{\}}
\end{alltt}
\end{kframe}
\end{knitrout}

\noindent Rebuild your package and check that the \cc{add} function still works.

\item Notice that we haven't exported the \cc{check\_numeric}
  function\sidenote{Remember to export a function, the function name should be
    in the \cc{NAMESPACE} file.}, so this will raise an error

\begin{knitrout}
\definecolor{shadecolor}{rgb}{0.969, 0.969, 0.969}\color{fgcolor}\begin{kframe}
\begin{alltt}
\hlkwd{library}\hlstd{(}\hlstr{"pkg"}\hlstd{)}
\hlkwd{check_numeric}\hlstd{(}\hlnum{1}\hlstd{)}
\end{alltt}
\end{kframe}
\end{knitrout}

\noindent but we can access any non-exported function in a package using the
\cc{:::} operator

\begin{knitrout}
\definecolor{shadecolor}{rgb}{0.969, 0.969, 0.969}\color{fgcolor}\begin{kframe}
\begin{alltt}
\hlstd{pkg}\hlopt{:::}\hlkwd{check_numeric}\hlstd{(}\hlnum{1}\hlstd{)}
\end{alltt}
\end{kframe}
\end{knitrout}

\noindent We can access\sidenote{The benefit of doing this is that we haven't
  loaded the package.} any exported function using \cc{::}
\begin{knitrout}
\definecolor{shadecolor}{rgb}{0.969, 0.969, 0.969}\color{fgcolor}\begin{kframe}
\begin{alltt}
\hlstd{pkg}\hlopt{::}\hlkwd{add}\hlstd{(}\hlnum{1}\hlstd{,} \hlnum{1}\hlstd{)}
\end{alltt}
\end{kframe}
\end{knitrout}
\item Now create a function \cc{subtract} and export this function. Rebuild your
  package and check that this works OK.
\item Delete your package and re-add the functions: \cc{add},
  \cc{check\_numeric} and \cc{subtract}\sidenote{The purpose of this is to
    highlight how easy it is to create packages.}.
\end{enumerate}


\newpage
\section{Documentation}

Using \cc{roxygen2} simplifies documentation\sidenote{The first package I wrote
  didn't use \cc{roxygen2} and it was a very painful experience.}. The premise of
\cc{roxygen2} is simple: describe your functions in comments next to their
definitions and \cc{roxygen2} will process your source code and comments to
produce \cc{Rd} files in the \cc{man/} directory. In theory, you should never
directly edit the \cc{Rd} files.

\begin{table}[b]
  \centering
  \begin{tabular}{@{} ll}
    \toprule
    Tag name & Description \\
    \midrule
    \cc{@title} & Short title for documentation page. \\
    \cc{@description} & Longer description page. Skip a line for a \\
    &\qquad new paragraph.\\
    \cc{@param} & Function parameter description. \\
    \cc{@inheritParams} & Use the parameter definition from another function. \\
    \cc{@export} & Add the function to the \cc{NAMESPACE} file.\\
    \cc{@return} & What does the function return, e.g. a data frame. \\
    \cc{@examples} & Function examples (will be run when building). \\
    \cc{@rdname} & Point multiple functions to the same help file,\\
    & \qquad e.g. \cc{?substr}. \\
    \cc{@seealso} & Pointers to other documentation pages. \\
    \cc{@importFrom} & Import functions from other packages. \\
    \bottomrule
  \end{tabular}
  \caption{Useful \cc{roxygen2} tags for documenting functions.}\label{T1.1}
\end{table}



\subsection{Tasks}

\begin{enumerate}
\item Copy the following \cc{roxygen2} descriptions to your add function
\begin{knitrout}
\definecolor{shadecolor}{rgb}{0.969, 0.969, 0.969}\color{fgcolor}\begin{kframe}
\begin{alltt}
\hlcom{#' @title A function for adding}
\hlcom{#'}
\hlcom{#' @description A really good adding function.}
\hlcom{#' Perhaps the best function ever!}
\hlcom{#'}
\hlcom{#' A work of pure genius.}
\hlcom{#' @param x a number}
\hlcom{#' @param y another number}
\hlcom{#' @return a number}
\hlcom{#' @export}
\hlcom{#' @examples}
\hlcom{#' add(5, 10)}
\hlcom{#' ## Can also use negative numbers}
\hlcom{#' add(-5, 10)}
\hlstd{add} \hlkwb{=} \hlkwa{function}\hlstd{(}\hlkwc{x}\hlstd{,} \hlkwc{y}\hlstd{) \{}
  \hlkwa{if}\hlstd{(}\hlopt{!}\hlkwd{check_numeric}\hlstd{(}\hlkwd{c}\hlstd{(x, y)))} \hlkwd{stop}\hlstd{(}\hlstr{"Not numeric"}\hlstd{)}
  \hlstd{x} \hlopt{+} \hlstd{y}
\hlstd{\}}
\end{alltt}
\end{kframe}
\end{knitrout}
\item Rebuild your package and look at the help file for the \cc{add} function,
  i.e. \cc{?add}. Run the examples via
\begin{knitrout}
\definecolor{shadecolor}{rgb}{0.969, 0.969, 0.969}\color{fgcolor}\begin{kframe}
\begin{alltt}
\hlkwd{example}\hlstd{(add)}
\end{alltt}
\end{kframe}
\end{knitrout}
\item Add a help page for the \cc{subtract} function.\sidenote{Build \& reload.}
\item Create a function called \cc{multiply} and add an associated help page.\sidenote{Build \& reload.}
\item Create a function called \cc{times}
\begin{knitrout}
\definecolor{shadecolor}{rgb}{0.969, 0.969, 0.969}\color{fgcolor}\begin{kframe}
\begin{alltt}
\hlstd{times} \hlkwb{=} \hlkwa{function}\hlstd{(}\hlkwc{x}\hlstd{,} \hlkwc{y}\hlstd{)} \hlkwd{multiply}\hlstd{(x, y)}
\end{alltt}
\end{kframe}
\end{knitrout}
\noindent Export the \cc{times} function.\sidenote{Build \& reload.}
\item Use the \cc{@rdname} tag above the \cc{times} function to point to the \cc{multiply} help page, i.e.
\begin{knitrout}
\definecolor{shadecolor}{rgb}{0.969, 0.969, 0.969}\color{fgcolor}\begin{kframe}
\begin{alltt}
\hlcom{#' @rdname multiply}
\end{alltt}
\end{kframe}
\end{knitrout}
\noindent Build and reload. Look at \cc{?multiply} and \cc{?times}. Now add
\cc{@examples} to the \cc{times} function. Look at the new \cc{times} help page.
\end{enumerate}

\subsection{Importing functions}

We often want to use functions from other R packages. When we do this we need to
be explicit, i.e. state what we want and from where. The great thing about R
packages, is that when we install a package, the dependencies are also
automatically installed.

\subsection*{Tasks}

\begin{enumerate}
\item Install the package \cc{jrAdvPackage}. First, we need the \cc{drat} package
\begin{knitrout}
\definecolor{shadecolor}{rgb}{0.969, 0.969, 0.969}\color{fgcolor}\begin{kframe}
\begin{alltt}
\hlkwd{install.packages}\hlstd{(}\hlstr{"drat"}\hlstd{)}
\end{alltt}
\end{kframe}
\end{knitrout}
\noindent Then we add the \cc{jr-packages} repo\sidenote{Run the \cc{.libPaths()}
  function to see the repository location.}
\begin{knitrout}
\definecolor{shadecolor}{rgb}{0.969, 0.969, 0.969}\color{fgcolor}\begin{kframe}
\begin{alltt}
\hlstd{drat}\hlopt{::}\hlkwd{addRepo}\hlstd{(}\hlstr{"jr-packages"}\hlstd{)}
\end{alltt}
\end{kframe}
\end{knitrout}
\noindent Then install as usual
\begin{knitrout}
\definecolor{shadecolor}{rgb}{0.969, 0.969, 0.969}\color{fgcolor}\begin{kframe}
\begin{alltt}
\hlkwd{install.packages}\hlstd{(}\hlstr{"jrAdvPackage"}\hlstd{)}
\end{alltt}
\end{kframe}
\end{knitrout}

\item The \cc{jrAdvPackage} contains a very useful function called \cc{div}
  that we want to use

\begin{knitrout}
\definecolor{shadecolor}{rgb}{0.969, 0.969, 0.969}\color{fgcolor}\begin{kframe}
\begin{alltt}
\hlkwd{library}\hlstd{(}\hlstr{"jrAdvPackage"}\hlstd{)}
\hlkwd{div}\hlstd{(}\hlnum{10}\hlstd{,} \hlnum{2}\hlstd{)}
\end{alltt}
\begin{verbatim}
## [1] 5
\end{verbatim}
\end{kframe}
\end{knitrout}

\noindent To use the \cc{div} function within our package, we have to import it first

\begin{knitrout}
\definecolor{shadecolor}{rgb}{0.969, 0.969, 0.969}\color{fgcolor}\begin{kframe}
\begin{alltt}
\hlcom{#' @importFrom jrAdvPackage div}
\end{alltt}
\end{kframe}
\end{knitrout}

\noindent and add an entry to the \cc{DESCRIPTION} file
\begin{verbatim}
    Imports: jrAdvPackage
\end{verbatim}
Create a function \cc{divide} that uses the \cc{div} function.
\end{enumerate}


\subsection{The \cc{DESCRIPTION} file}

The \cc{DESCRIPTION} file contains high level information about your package.
For example, the package name, a brief description, the licence, and your email
address.

Open the \cc{DESCRIPTION} file and update fields with relevant information. An
example is given below.
\begin{verbatim}
Package: pkg
Type: Package
Title: My First package
Version: 0.1
Date: 2016-11-01
Authors@R: person(given="Colin", family="Gillespie",
    email="colin@jumpingrivers.com", role = c("aut", "cre"))
Maintainer: Colin Gillespie <colin@jumpingrivers.com>
Description: This is my very first package. It contains
  exceedingly useful functions, such as add and subtract.
  Make sure you add a couple of spaces to indent the
  Description otherwise you will waste hours of your life
  trying to find the bug.
License: GPL-2 | GPL-3
LazyData: TRUE
\end{verbatim}

\subsection{Package checks}

One of great things about R packages, is that there are a number of package
checks that are available. These include
\begin{itemize}
\item Checking the syntax of the \cc{DESCRIPTION} and \cc{NAMESPACE} file.
\item Checking your examples run.
\item Checking all exported functions are documented.
\end{itemize}

\subsection*{Tasks}

\begin{enumerate}
\item Run the standard package checks on your package, via
\begin{verbatim}
    Build -> Check Package
\end{verbatim}
  Check that you package passes all tests. \sidenote{CTRL + E} Fix any errors,
  warnings or notes.
\item Add the following example to the \cc{add} function

\begin{knitrout}
\definecolor{shadecolor}{rgb}{0.969, 0.969, 0.969}\color{fgcolor}\begin{kframe}
\begin{alltt}
\hlcom{#' add("A", "B")}
\end{alltt}
\end{kframe}
\end{knitrout}

\noindent Build the package. Does the package still build? Check the package.
Does the package pass all tests?
\end{enumerate}

\newpage
\section{Data and demos}

\subsection{Data in packages}

Packages can also contain example data sets\sidenote{Full details are given at
  \mbox{\url{http://goo.gl/Y4Srx4}}.}. Data files\sidenote{Note that R code
  should be "self-sufficient" and not make use of extra functionality provided
  by the package, so that the data file can also be used without having to load
  the package.} can be one of three types as indicated by their extension.
\begin{itemize}
\item plain R code (\cc{.R} or \cc{.r})
\item tables (\cc{.tab}, \cc{.txt}, or \cc{.csv})
\item \cc{save()} images (\cc{.RData} or \cc{.rda}).\marginnote{I tend to use
    \cc{RData files.}}
\end{itemize}
Data files live in the \cc{data/} directory.

Each data file should also have an associated help page. The easiest way to
generate a help page is to use \cc{roxygen2} and a dummy R function. Typically,
I have a file called \cc{data\_help\_files.R}, which has entries for the each
data set. For example,\marginnote{This is entry is taken from the \cc{poweRlaw}
  package.}
\begin{knitrout}
\definecolor{shadecolor}{rgb}{0.969, 0.969, 0.969}\color{fgcolor}\begin{kframe}
\begin{alltt}
\hlcom{## Entry in data_help_files.R}
\hlcom{## Name is name of the data set.}
\hlcom{#' @name moby}
\hlcom{#' @aliases moby_sample}
\hlcom{#' @title Moby Dick word count}
\hlcom{#' @description The frequency of occurrence of unique words}
\hlcom{#' in the novel Moby Dick by Herman Melville.}
\hlcom{#'}
\hlcom{#' The data set moby_sample is 2000 values sampled from the}
\hlcom{#' moby data set.}
\hlcom{#' @docType data}
\hlcom{#' @format A vector}
\hlcom{#' @source M. E. J. Newman, "Power laws, Pareto distributions}
\hlcom{#' and Zipf's law." Contemporary Physics 46, 323 (2005).}
\hlkwa{NULL}
\end{alltt}
\end{kframe}
\end{knitrout}

\subsection{Tasks}

\begin{enumerate}
\item Create a \cc{data/} directory.
\item Create the following data frame
\begin{knitrout}
\definecolor{shadecolor}{rgb}{0.969, 0.969, 0.969}\color{fgcolor}\begin{kframe}
\begin{alltt}
\hlstd{example_data} \hlkwb{=} \hlkwd{data.frame}\hlstd{(}\hlkwc{x} \hlstd{=} \hlkwd{runif}\hlstd{(}\hlnum{10}\hlstd{),} \hlkwc{y} \hlstd{=} \hlkwd{runif}\hlstd{(}\hlnum{10}\hlstd{))}
\end{alltt}
\end{kframe}
\end{knitrout}
\noindent Now save\sidenote{Use the \cc{save} function.} the data frame \cc{example\_data}
in the \cc{data/} directory
\begin{knitrout}
\definecolor{shadecolor}{rgb}{0.969, 0.969, 0.969}\color{fgcolor}\begin{kframe}
\begin{alltt}
\hlkwd{save}\hlstd{(example_data,} \hlkwc{file}\hlstd{=}\hlstr{"data/example_data.RData"}\hlstd{)}
\end{alltt}
\end{kframe}
\end{knitrout}
\item Create a file called \cc{data\_help\_files.R} in the \cc{R/} directory and
  document your new data set.
\item Build and reload your package. Can you access the help page and the data
  set?
\item Check that your package still passes all tests\sidenote{CTRL + E}.
\end{enumerate}


\subsection{Demos}

A demo is similar to function examples, but is typically longer and shows how to
use multiple functions together. Demos are plain \cc{.R} files that live in the
\cc{demo/} package directory. The demos are accessed with the \cc{demo()}
function.

In \cc{demo/} directory, there should also be an \cc{00Index} file, that lists
the demos\sidenote{There is a planned \cc{demoTitle} tag for \cc{roxygen2}, but
  currently this hasn't been implemented.}. For example,\sidenote{Note the white
  space separation in the \cc{00Index}. Use at least four spaces to avoid
  annoying error messages.}
\begin{verbatim}
demo1      My very first demo
demo2      My very second demo
\end{verbatim}

\subsection{Tasks}

\begin{enumerate}
\item Create a \cc{demo/} directory.
\item Create a file called \cc{first.R} and save it in the \cc{demo/} directory.
  In this file show how you can use some of your newly created function.
\item Add a \cc{00Index} file to the \cc{demo/} directory.
\item Build and check your package.
\end{enumerate}

\newpage
\section{Vignettes}

If you want to include more extensive examples or even just further
documentation, then you should consider creating a vignette:
\begin{quote}
  a \textbf{vignette} is a small illustration placed at the beginning or end of
  a book or chapter.

  \em \url{http://dictionary.reference.com/browse/vignette}
\end{quote}
We can view vignettes from other packages using the \cc{vignette} function

\begin{knitrout}
\definecolor{shadecolor}{rgb}{0.969, 0.969, 0.969}\color{fgcolor}\begin{kframe}
\begin{alltt}
\hlkwd{vignette}\hlstd{(}\hlkwc{package}\hlstd{=}\hlstr{"knitr"}\hlstd{)}
\end{alltt}
\end{kframe}
\end{knitrout}
\noindent or to view in your web browser

\begin{knitrout}
\definecolor{shadecolor}{rgb}{0.969, 0.969, 0.969}\color{fgcolor}\begin{kframe}
\begin{alltt}
\hlkwd{browseVignettes}\hlstd{(}\hlkwc{package}\hlstd{=}\hlstr{"knitr"}\hlstd{)}
\end{alltt}
\end{kframe}
\end{knitrout}
\subsection{Example: markdown vignettes}

Vignettes are stored in the \cc{vignettes/} directory. The simplest vignette
uses R-markdown and is formatted by the \cc{knitr} package. To create a package
vignette, we simply place the file in the \cc{vignettes/} directory.
\marginnote{The output style is \cc{html\_vignette}. This is more lightweight (in terms
of file size) than the standard \cc{html\_document}} For
example, suppose we have a file \cc{intro.Rmd} that contains the following text:
\marginnote{The \cc{UTF-8} line specifies the file encoding.}
\begin{knitrout}
\definecolor{shadecolor}{rgb}{0.969, 0.969, 0.969}\color{fgcolor}\begin{kframe}
\begin{alltt}
---
title: \hlstr{"My very first vignette"}
author: \hlstr{"Colin Gillespie and Robin Lovelace"}
output: rmarkdown::html_vignette
vignette: >
  %\textbackslash{}VignetteIndexEntry\{My very first vignette\}
  %\textbackslash{}VignetteEngine\{knitr::rmarkdown\}
  %\textbackslash{}VignetteEncoding\{UTF-8\}
---

\hlcom{## My first package}

This is my **first** package vignette.
I can include mathematics, such as $x^2$.
R code is also nicely formatted and displayed.

```\{r\}
x = \hlkwd{runif}(10)
x
```
and plots
```\{r\}
\hlkwd{plot}(x)
```
\end{alltt}
\end{kframe}
\end{knitrout}

\begin{enumerate}
\item The first few lines register \cc{knitr} as the vignette engine and provide
  an entry for the list of vignettes.
\item We have used markdown to add simple styling\sidenote{See
    \url{http://goo.gl/y9FO5} and the RStudio markdown reference sheet for
    further commands.}. For example \cc{**first**} becomes \textbf{first} and
  \cc{\$x\textasciicircum 2\$} becomes $x^2$.
\item R code is executed in the ``` regions.
\end{enumerate}

\subsection{Tasks}

\begin{enumerate}
\item Create a \cc{vignettes/} directory in your package.
\item Create an R markdown file
\begin{verbatim}
    File -> New File -> R markdown
\end{verbatim}
and save it in the \cc{vignettes/} directory.
\item Copy the markdown example above into your file and \textit{knit} that file.
\item Add
\begin{verbatim}
    Suggests: knitr
    VignetteBuilder: knitr
\end{verbatim}
  to the \cc{DESCRIPTION} file\sidenote{Notice the \cc{knitr} is only a
    \textbf{suggested} package, since it's not essential.}. This will tell
  your package to build vignettes using \cc{knitr}.
\item Vignettes won't actually be built unless you are creating a source
  bundle\sidenote{If you submit your package to CRAN, then the vignettes will
    be built ready for distribution.}. To install a package with vignettes included, we
  first create the source package
\begin{verbatim}
    Build -> Build Source Package
\end{verbatim}
and then install the package from source
\begin{knitrout}
\definecolor{shadecolor}{rgb}{0.969, 0.969, 0.969}\color{fgcolor}\begin{kframe}
\begin{alltt}
\hlkwd{install.packages}\hlstd{(}\hlstr{"pkg_1.0.tar.gz"}\hlstd{,} \hlkwc{repos}\hlstd{=}\hlkwa{NULL}\hlstd{,}
                 \hlkwc{type}\hlstd{=}\hlstr{"source"}\hlstd{)}
\end{alltt}
\end{kframe}
\end{knitrout}
\noindent Build and install your package. Check that you can access the
vignette.
\item Check and ensure that your package passes all tests\sidenote{CTRL + E}.
\end{enumerate}

\subsection{Package level documentation}

You package typically also has documentation associated with the package name.
For example,
\begin{knitrout}
\definecolor{shadecolor}{rgb}{0.969, 0.969, 0.969}\color{fgcolor}\begin{kframe}
\begin{alltt}
\hlkwd{library}\hlstd{(}\hlstr{"pkg"}\hlstd{)}
\hlopt{?}\hlstd{pkg}
\end{alltt}
\end{kframe}
\end{knitrout}
\noindent should bring up an overview of your package. Again we use \cc{roxygen2}
to generate the man page. Go to
\begin{center}
  \url{http://goo.gl/W2tJrF}
\end{center}
to view the entry for the \cc{poweRlaw} package.

\subsection{Tasks}

\begin{enumerate}
\item Create a file called \cc{pkg-package.R} in the \cc{R/} directory.
\item Using the \cc{poweRlaw} package as an example, create a man page for your
  package\sidenote{What do you think the \cc{@aliases} tag does?}.
\item Do one final check and ensure that your package passes all tests.
\end{enumerate}



\end{document}
